\documentclass{article}
\usepackage{/globals/de}
\begin{document}
\section{Betriebsysteme} 
Ein Betriebssystem (\emph{Operating System}, \emph{OS}) ist das erste Programm das auf einem Computer läuft.
Es übernimmt die Verwaltung von 
\begin{description}
 \item[Die Hardware] Über vom Betriebssystem zur Verfügung gestellte Treiber kann der Computer mit Ein- und Ausgabegeräten interagieren. 
 \item[Das Dateisystem] Dateien werden über das Betriebssystem gefunden, gelesen, geschrieben und beschrieben.
 \item[Benutzeroberflächen] Das Betriebssystem zeigt ist dafür verantwortlich den Desktop, alle offenen Programme, extras wie \zb eine Taskleiste, \etc anzuzeigen.
 \item[BenutzerInnen] Ein Computer kann oftmals von mehreren BenutzerInnen, mit seperaten Dateien und Rechten, genutzt werden.
 \item[Prozesse] Das Betriebssystem entscheided welche Prozesse wie viel Zeit auf der CPU bekommen, welchen RAM sie nutzen dürfen, oder ob sie nicht vorzeitig beendet werden sollten wenn es nicht mehr genug Ressourcen gibt.
\end{description}

\end{document}